% year: תשפ"ב
% type: מבחן
% semester: א
% moed: מיוחד
% number: 2ב
% points: 9
% categories: ivt, mvt
% source: exams/מבחנים/תשפב/מועד ג תשפב.pdf

\begin{question}
הוכיחו כי לפונקציה $f(x)=\frac15x^5+\frac23x^3+x+7$ יש שורש ממשי יחיד.
\end{question}

\begin{hint}
הקיום נובע ממשפט ערך הביניים, והיחידות — ממונוטוניות: $f'(x)=x^4+2x^2+1$ (או ממשפט רול).
\end{hint}

\begin{solution}
\textbf{קיום.} $f$ פולינום ולכן רציפה בכל $\R$. מתקיים $f(0)=7>0$ ו-
\[ f(-2)=-\frac{32}{5}-\frac{16}{3}-2+7=-\frac{101}{15}<0. \]
לפי \textbf{משפט ערך הביניים} קיים $x_0\in(-2,0)$ עם $f(x_0)=0$.

\textbf{יחידות.}
\[ f'(x)=x^4+2x^2+1=(x^2+1)^2>0\quad\text{לכל } x, \]
ולכן $f$ עולה ממש ב-$\R$, ובפרט חד-חד-ערכית — היא מקבלת את הערך $0$ לכל היותר פעם אחת.
(לחלופין, לפי משפט רול: אם היו שני שורשים $x_1<x_2$, היה קיים $c\in(x_1,x_2)$ עם $f'(c)=0$, בסתירה ל-$f'>0$.)

לכן ל-$f$ יש שורש ממשי יחיד, והוא נמצא בקטע $(-2,0)$.
\end{solution}
